\documentclass{article}

\usepackage[utf8]{inputenc}
\usepackage[T1]{fontenc}
\usepackage[spanish]{babel} 
\usepackage{graphicx}       
\usepackage{amsmath}        
\usepackage{amssymb}
\usepackage{float}
\usepackage{draftwatermark} 
\bibliographystyle{ieeetr} 

\usepackage{verbatim}

\SetWatermarkText{EN REVISIÓN}
\SetWatermarkScale{3}
\SetWatermarkColor[gray]{0.85}


\title{CAT797 lateral dynamics}
%\author{Jiménez Ramírez}
\date{\today}

\begin{document}

\maketitle

Sea la dinámica lateral del vehículo descrita por el modelo bicicleta lineal

\begin{equation}
\ddot{y}
=
\left(
\frac{C_r l_r - C_f l_f}{m \dot{x}}
-
\dot{x}
\right)\dot{\psi}
-
\left(
\frac{C_f + C_r}{m \dot{x}}
\right)\dot{y}
+
\frac{C_f}{m}\delta_f
\label{eq:lateral_dynamics}
\end{equation}

\begin{equation}
\ddot{\psi}
=
-
\left(
\frac{C_f l_f^2 + C_r l_r^2}{I_z \dot{x}}
\right)\dot{\psi}
+
\left(
\frac{C_r l_r - C_f l_f}{I_z \dot{x}}
\right)\dot{y}
+
\frac{C_f l_f}{I_z}\delta_f
\label{eq:yaw_dynamics}
\end{equation}

\begin{equation}
\dot{Y}
=
\dot{x}\sin\psi+\dot{y}\cos\psi
\end{equation}

Suponiendo pequeños ángulos de guiñada, es posible aproximar

\begin{equation}
\dot{Y}
\approx
\dot{x}\psi+\dot{y}
\label{eq:Ydot}
\end{equation}

Definiendo el vector de estado

\begin{equation}
\chi=
\begin{bmatrix}
\dot{y}\\
\dot{\psi}\\
Y\\ \psi
\end{bmatrix}
=
\begin{bmatrix}
\chi_1\\
\chi_2\\
\chi_3\\ \chi_4
\end{bmatrix}
\end{equation}

y tomando como salida la posición lateral

\begin{equation}
y_\chi=
\begin{bmatrix}
0 & 0 & 1 & 0
\end{bmatrix}
\chi
=
\chi_3
\end{equation}

El error de seguimiento se define como

\begin{equation}
e=y_\chi-y_{ref}
\end{equation}

donde $y_{ref}$ representa la trayectoria lateral deseada.

Se selecciona una superficie deslizante de primer orden dada por

\begin{equation}
\sigma=\dot e+\lambda e
\label{eq:sliding_surface}
\end{equation}

donde $\lambda>0$ determina la velocidad de convergencia del error.

Derivando \eqref{eq:sliding_surface}

\begin{equation}
\dot{\sigma}
=
\ddot e+\lambda \dot e
\label{eq:sigma_dot}
\end{equation}

Por otra parte,

\begin{equation}
e=y_\chi-y_{ref}
\end{equation}

por lo que

\begin{equation}
\ddot e
=
\ddot y_\chi-\ddot y_{ref}
\label{eq:error_ddot}
\end{equation}

Dado que

\begin{equation}
y_\chi=Y
\end{equation}

de \eqref{eq:Ydot} se obtiene

\begin{equation}
\dot y_\chi
=
\dot{x}\psi+\dot y
= \dot{x}\chi_4+\chi_1
\end{equation}

Derivando nuevamente

\begin{equation}
\ddot y_\chi
=
\ddot{x}\psi
+
\dot{x}\dot{\psi}
+
\ddot y
= \ddot{x}\chi_4 + \dot{x}\chi_2 + \dot{\chi}_1
\label{eq:yddot_aux}
\end{equation}

Sustituyendo la dinámica lateral \eqref{eq:lateral_dynamics} en \eqref{eq:yddot_aux}

\begin{equation}
\begin{aligned}
\ddot y_\chi
=&
\ddot{x}\chi_4
+
\dot{x}\chi_2
+
\left(
\frac{C_r l_r-C_f l_f}{m\dot{x}}
-\dot{x}
\right)\chi_2
\\
&
-\frac{C_f+C_r}{m\dot{x}}\chi_1
+
\frac{C_f}{m}\delta_f
\end{aligned}
\end{equation}

Agrupando términos

\begin{equation}
\ddot y_\chi
=
\ddot{x}\chi_4
+
\frac{C_r l_r-C_f l_f}{m\dot{x}}\chi_2
-
\frac{C_f+C_r}{m\dot{x}}\chi_1
+
\frac{C_f}{m}\delta_f
\label{eq:input_output}
\end{equation}

Se define una entrada virtual $v$ como

\begin{equation}
v=\ddot y_\chi
\end{equation}

por lo que de \eqref{eq:input_output}

\begin{equation}
v
=
\ddot{x}\chi_4
+
\frac{C_r l_r-C_f l_f}{m\dot{x}}\chi_2
-
\frac{C_f+C_r}{m\dot{x}}\chi_1
+
\frac{C_f}{m}\delta_f
\end{equation}

Despejando la entrada de control

\begin{equation}
\delta_f
=
\frac{m}{C_f}
\left(
v
-\ddot{x}\chi_4
-\frac{C_r l_r-C_f l_f}{m\dot{x}}\chi_2
+\frac{C_f+C_r}{m\dot{x}}\chi_1
\right)
\label{eq:feedback_linearization}
\end{equation}

Para reducir el fenómeno de chattering se emplea una capa límite basada en la función hiperbólica tangente

\begin{equation}
\dot{\sigma}
=
-k\tanh\left(\frac{\sigma}{\nu_\chi}\right)
\label{eq:reaching_law}
\end{equation}

donde $k>0$ es la ganancia de convergencia y $\nu_\chi>0$ representa el espesor de la capa límite.

Sustituyendo \eqref{eq:error_ddot} en \eqref{eq:sigma_dot}

\begin{equation}
\dot{\sigma}
=
\ddot y_\chi
-\ddot y_{ref}
+\lambda \dot e
\end{equation}

y utilizando \eqref{eq:reaching_law}

\begin{equation}
\ddot y_\chi
=
\ddot y_{ref}
-\lambda\dot e
-k\tanh\left(\frac{\sigma}{\nu_{sw}}\right)
\end{equation}

Por lo tanto, la entrada virtual queda definida como

\begin{equation}
v
=
\ddot y_{ref}
-\lambda\dot e
-k\tanh\left(\frac{\sigma}{\nu_{sw}}\right)
\label{eq:v}
\end{equation}

Sustituyendo \eqref{eq:v} en \eqref{eq:feedback_linearization}, se obtiene la ley de control final en términos de $\chi$

\begin{equation}
\boxed{
\delta_f=
\frac{m}{C_f}
\left(
\ddot y_{ref}
-\lambda\dot e
-k\tanh\left(\frac{\sigma}{\nu_{sw}}\right)
-\ddot{x}\chi_4
-\frac{C_r l_r-C_f l_f}{m\dot{x}}\chi_2
+\frac{C_f+C_r}{m\dot{x}}\chi_1
\right)
}
\label{eq:control_law}
\end{equation}

en este caso si $\dot{x}=$constante, $\ddot{x}=0 $, la ley de control se simplifica a

\begin{equation}
\boxed{
\delta_f=
\frac{m}{C_f}
\left(
\ddot y_{ref}
-\lambda\dot e
-k\tanh\left(\frac{\sigma}{\nu_{sw}}\right)
-\frac{C_r l_r-C_f l_f}{m\dot{x}}\chi_2
+\frac{C_f+C_r}{m\dot{x}}\chi_1
\right)
}
\label{eq:control_law_simplified}
\end{equation}

\end{document}
